% =========================================================================
% Capítulo 4: Cuerpo rígido
% Tema 01: Cuerpo rígido
% Problema 09: Roca esférica empujada por una topadora cuesta arriba
% =========================================================================
\subsubsection{Roca esférica empujada por una topadora cuesta arriba}

Una topadora empuja una roca de masa $m$ con una fuerza conocida $\mathbf{P}$ cuesta arriba
por una ladera inclinada un ángulo constante $\beta$, como se muestra en la Figura
\ref{fig:c4_e01_p09_topadora}. La roca se modela como una esfera uniforme de masa $m$ y radio
$r$. Se sabe que la roca rueda sin deslizar sobre la ladera, que el coeficiente de fricción
dinámica de Coulomb entre la hoja de la topadora y la roca es $\mu$, que la fuerza
$\mathbf{P}$ es paralela a la pendiente y pasa por el centro de la roca, y que la gravedad
actúa verticalmente hacia abajo.

\begin{enumerate}
    \item Determine la ecuación diferencial de movimiento de la roca en términos del
          desplazamiento $x$ a lo largo de la pendiente.
\end{enumerate}

\textbf{Pista:} el momento de inercia de una esfera maciza uniforme respecto a un eje por su
centro es $I=\tfrac25 mr^2$.

\begin{figure}[!ht]
    \centering
    \begin{tikzpicture}[>=Latex, scale=1.1]

        % --- PARAMETROS ---
        \def\B{28}       % Angulo de la ladera
        \def\r{0.7}      % Radio de la roca

        % --- SUELO HORIZONTAL ---
        \draw[thick] (-0.3,0) -- (7.5,0);

        % --- LADERA (rotada) ---
        \begin{scope}[rotate=\B]
            % superficie inclinada
            \draw[thick] (0,0) -- (7.2,0);
            \foreach \x in {0.3,0.7,...,7.0} { \draw (\x,0) -- ++(-0.18,-0.18); }

            % roca (esfera)
            \coordinate (O) at (3.4,\r);
            \draw[thick, fill=gray!20] (O) circle (\r);
            \filldraw (O) circle (1pt) node[above right=-1pt] {$O$};
            \draw (O) -- ++(-40:\r) node[midway, below left=-3pt] {$r$};

            % hoja de la topadora (placa) del lado cuesta abajo
            \draw[line width=2.5pt, gris1] ({3.4-\r-0.05},-0.05) -- ({3.4-\r-0.05},{2*\r});
            \node[gris1] at ({3.4-\r-0.6},{1.5*\r}) {topadora};

            % fuerza P paralela a la pendiente por el centro
            \draw[->, very thick, myred] ($(O)+(-1.6,0)$) -- ($(O)+(-0.75,0)$) node[midway, above] {$\mathbf{P}$};

            % coordenada x a lo largo de la pendiente
            \draw[->] (4.6,-0.9) -- (5.6,-0.9) node[right] {$x$};

            % contacto sin deslizar
            \node[font=\small] at (5.2,0.5) {no desliza};
            \draw (4.9,0.35) -- (3.9,0.02);

            % friccion en la hoja
            \node[font=\small] at (2.0,{2.2*\r}) {fricción $\mu$};
        \end{scope}

        % --- ANGULO BETA ---
        \draw[->] (1.6,0) arc (0:\B:1.6);
        \node at ({\B/2}:1.85) {$\beta$};

        % --- GRAVEDAD ---
        \draw[->, thick] (7.0,2.2) -- (7.0,1.4) node[right, midway] {$\mathbf{g}$};

    \end{tikzpicture}
    \caption{Roca esférica ($m,r$) empujada con fuerza $\mathbf{P}$ paralela a una ladera de inclinación $\beta$; rueda sin deslizar sobre la ladera.}
    \label{fig:c4_e01_p09_topadora}
\end{figure}

% --- SOLUCIÓN ---------------------------------------------------------------
\begin{solucion}{Roca esférica empujada por una topadora cuesta arriba}
\subsubsection*{Marcos y cinemática}

Se usa un marco inercial fijo a la ladera con $\vec{E}_x$ cuesta arriba (a lo largo de la
pendiente), $\vec{E}_y$ normal a la superficie y $\vec{E}_z=\vec{E}_x\times\vec{E}_y$. El
centro de la esfera está en $\vec{r}_O=x\vec{E}_x$, de modo que $\vec{a}_O=\ddot{x}\vec{E}_x$.
La condición de rodadura sin deslizar en el punto de contacto $Q$ con la ladera da
$$\dot{x}=r\omega \;\Longrightarrow\; \omega=\frac{\dot{x}}{r},\qquad \dot\omega=\frac{\ddot{x}}{r}.$$

\subsubsection*{Fuerzas}

El punto de la esfera en contacto con la hoja ($P$) desliza respecto a ésta, por lo que la
fricción dinámica de la hoja es $\vec{F}_f=\mu P\,\vec{E}_y$ (magnitud $\mu P$). Además actúan
la fuerza $\vec{P}=P\vec{E}_x$, el peso $\vec{W}=-mg\sin\beta\,\vec{E}_x + mg\cos\beta\,\vec{E}_y$,
y la normal y la tracción del suelo en $Q$ (que no generan torque respecto a $Q$).

\subsubsection*{Balance de momento angular respecto al contacto $Q$}

Para evitar las reacciones desconocidas del suelo, se toma momento angular respecto al punto
de contacto $Q$. Con $\bar{I}=\tfrac25 mr^2$:
$$\vec{H}_Q = \bar{I}\,\omega\,\vec{E}_z + (\vec{r}_O-\vec{r}_Q)\times m\vec{v}_O
= \frac{2}{5}mr\dot{x}\,\vec{E}_z + mr\dot{x}\,\vec{E}_z = \frac{7}{5}mr\dot{x}\,\vec{E}_z,$$
de donde $\dot{\vec{H}}_Q = \tfrac75 mr\ddot{x}\,\vec{E}_z$.

El torque respecto a $Q$, con $\vec{r}_P-\vec{r}_Q=-r\vec{E}_x-r\vec{E}_y$ y
$\vec{r}_O-\vec{r}_Q=-r\vec{E}_y$, resulta
$$\vec{M}_Q = (\vec{r}_P-\vec{r}_Q)\times(\vec{P}+\vec{F}_f) + (\vec{r}_O-\vec{r}_Q)\times\vec{W}
= \big[rP(1-\mu) - mgr\sin\beta\big]\vec{E}_z.$$

Igualando $\vec{M}_Q=\dot{\vec{H}}_Q$ y dividiendo por $r$:
$$rP(1-\mu) - mgr\sin\beta = \frac{7}{5}mr\ddot{x}
\;\Longrightarrow\;
\boxed{\;\frac{7}{5}m\ddot{x} + mg\sin\beta = P(1-\mu)\;}$$

Equivalentemente, $\ddot{x} = \dfrac{5}{7m}\big[P(1-\mu) - mg\sin\beta\big]$.
\end{solucion}
